% p.60 — الباب السادس والعشرون في المصادر والصفة المشبهة واسم المبالغة
% Layout (as in the lithograph): one boxed table of six narrow columns of word lists
% (21 cells high), read top to bottom, right to left, with the prose set INSIDE the grid
% where it stands on the page: (a) the top of column 1 (rows 1–7) is one prose cell (rule on
% مصادر الثلاثي المجرد), set as a \multirow over seven rows; column 1 then holds مَصْدَرُ فَتَحَ …
% (14 cells); (b) column 4, under the sub-heading «صفة مشبهة» (row 8), holds the definition
% (rows 9–19, one \multirow cell) followed by the two examples أَحْمَرُ / حَمْرَاءُ;
% (c) column 6 has the sub-heading «اسم مبالغة» with a two-line definition, kept as two cells.
% Sub-headings inside columns (مَصْدَرُ فَتَحَ etc.) are \thead cells. Columns 1 and 4 are
% p-columns (centred) so that the prose cells wrap.
\mspage{60}
\nohairlines
\bab{الباب السادس والعشرون في المصادر والصفة المشبهة واسم المبالغة}
\begin{sarf}{|>{\centering\arraybackslash}p{30mm}|c|c|>{\centering\arraybackslash}p{30mm}|c|c|}
\multirow{7}{30mm}{اعلم ان مصادر الثلاثي المجرد غير قياسي لا يكون على نهج واحد بخلاف مصادر الثلاثي المزيد و الرباعي فانها على نهج واحد قياسي.} & مَكْذُوْبَةً & تَجْوَالًا & \thead{مَصْدَرُ كَرُمَ} & حُمْرٌ & غَيُوْرٌ\\
 & \unclear{حَيُّوْرَةً} & \unclear{جَبْرُوَّةً} & كَرَامَةً & أَعْرَجُ & حُبْلٰى\\
 & \unclear{عُلْبٰى} & كَيْنُوْنَةً & صِغَرًا & عَرْجَاءُ & حَيْوَانٌ\\
 & \unclear{تَيْسِيْرًا} & دِلِّيْلٰى & كَثْرَةً & عُرْجٌ & صِفْرٌ\\
 & غُفْرَانًا & \thead{مَصْدَرُ سَمِعَ} & قُرْبًا & صَعْبٌ & \thead{اسم مبالغة}\\
 & مَكْذُوْبًا & رَحْمَةً & قُرْبَانًا & صُلْبٌ & ما فيه معنى\\
 & كَاذِبَةً & كُدْرَةً & وِلَايَةً & حَسَنٌ & وصف مع زيادة\\\cline{1-1}
\thead{مَصْدَرُ فَتَحَ} & مَمْلَكَةً & \unclear{دَهَارَةً} & \thead{صفة مشبهة} & خَشِنٌ & سَمَّاكٌ\\
قَطِيْعَةً & \thead{مَصْدَرُ نَصَرَ} & \unclear{مَحْمُوْدَةً} & \multirow{11}{30mm}{ما يدل على وصف ثابت في شئ وان كان فيه معنى لون او عيب يكون على وزن افعل مثل احمر و اعرج ويكون المؤنث فيه على وزن فعلاء مثل حمراء و عرجاء ويكون الجمع على وزن فُعْلٍ.} & نَدُسٌ & كُبَّارٌ\\
ذَهَابًا & قَتْلًا & قَبُوْلًا & & \unclear{دِيَمٌ} & صِدِّيْقٌ\\
\unclear{يَقْطَاعًا} & فِسْقًا & \unclear{رَغْبَاءً} & & \unclear{بِلِزٌ} & قَيُّوْمٌ\\
مَشْعَاةً & شُكْرًا & \unclear{دَغْبُوَةً} & & جُنُبٌ & قُدُّوْسٌ\\
\thead{مَصْدَرُ ضَرَبَ} & نِشْدَةً & مَحْمِدَةً & & عَطْشَانُ & عُجَابٌ\\
\unclear{خَنِيْقًا} & طَلَبًا & كَرَاهِيَةً & & عُرْيَانُ & فَارُوْقٌ\\
دِرَايَةً & قِيَامًا & \unclear{رَغْبُوْلٰى} & & كَابِرٌ & هُمَزَةٌ\\
\unclear{بُحَايَةً} & نَزَوَانًا & \thead{مَصْدَرُ حَسِبَ} & & سَيِّدٌ & حَذِرٌ\\
غَلَبَةً & مَدْخَلًا & حُسْبَانًا & & \unclear{جَيَّانٌ} & عَلِيْمٌ\\
سَرِقَةً & دُخُوْلًا & نِعْمَةً & & هِجَانٌ & أَكُوْلٌ\\
هُدًى & دَعْوٰى & وَرَعًا & & وَضَّاعٌ & \unclear{مِحْرَمٌ}\\
قَيْلُوْلَةً & بُشْرٰى & وِلَايَةً & أَحْمَرُ & كُبَّارٌ & مِنْعَامٌ\\
لَيَّانًا & ذِكْرٰى & يَأْسًا & حَمْرَاءُ & كَرِيْمٌ & مِنْطِيْقٌ\\
\end{sarf}
